\chapter{Projeto Conceitual do Produto}

\section{Características gerais}

\textcolor{red}{Descrever produto de forma geral. Apresentar itens teóricos sobre o projeto a serem aprofundados ou detalhados oportunamente.}

\textcolor{red}{Apresentar e explicar a Estrutura Analítica do Projeto (EAP), tomando cuidado com a legibilidade da imagem. Se necessário, coloque a EAP dentro do comando ``\textsf{\textbackslash begin\{landscape\} \textbackslash end\{landscape\}}'' para que ela seja apresentada em uma página deitada.}

\section{Estrutura}

\textcolor{red}{Apresentar desenho em CAD, indicar dimensões por lado/aresta/cota, indicar materiais utilizados e explicar o desenho e as decisões de projeto.}

\section{Descrição de \textit{hardware}}

\textcolor{red}{A descrição de \textit{hardware} no documento deve permitir que outras pessoas repliquem o produto proposto neste documento. Para isso, deve-se explicar de forma textual os principais sub-sistemas e as escolhas de projeto, além de indicar:}

\begin{itemize}
    \item \textcolor{red}{ O diagrama de blocos \cite{blockdiagram}, que oferece uma visão geral do hardware, com as principais partes do sistema e as ligações entre elas.}
    %\item A lista de materiais \cite{bom}, que facilita a montagem do sistema, indicando tudo que deve ser adquirido. Aproveitem a lista de materiais para apresentarem o orçamento dos mesmos.
    \item \textcolor{red}{ O esquemático \cite{esquematico}, que mostra as conexões elétricas entre os componentes. Deve-se indicar os componentes por nomes e símbolos, e as conexões entre eles, incluindo a pinagem dos componentes. Se o esquemático ficar muito grande, pode-se separa-lo em vários esquemáticos. A ligação em \textit{protoboard} não é considerada um esquemático adequado.}
\end{itemize}

\textcolor{red}{A Fig. \ref{fig:diagrama_blocos_e_esquematico} apresenta o diagrama de blocos e o esquemático de um circuito conversor de corrente alternada para corrente direta, a título de exemplo. Repare que o diagrama de blocos indica os subsistemas do circuito completo apresentado no esquemático.
% Comecem com o diagrama de blocos, que é mais genérico, e depois indiquem os detalhes com a lista de materiais e os esquemáticos.
}

\begin{figure}
    \centering
    \includegraphics[width=0.5\linewidth]{figuras/Exemplo-diagrama-de-blocos.png}
    \includegraphics[width=0.5\linewidth]{figuras/Esquematico.png}
    \caption{\textcolor{red}{Diagrama de blocos e esquemático de um circuito conversor de corrente alternada para corrente direta. Extraído de \cite{block_n_schematic}.}}
    \label{fig:diagrama_blocos_e_esquematico}
\end{figure}

\section{Análise de consumo energético}

\textcolor{red}{Com relação ao consumo energético do produto desenvolvido, será necessário apresentar cálculos de consumo dos diversos componentes utilizados e explicar as decisões de projeto para atender a estas demandas.}


\begin{itemize}
    \item \textcolor{red}{ \textbf{Identificação dos Subsistemas Elétricos:} liste todos os componentes que consumirão energia, como sensores, microcontroladores, atuadores e módulos de armazenamento de dados. Para cada componente, registre a tensão de operação (V), corrente média (mA) e tempo estimado de uso (s).}
    \item \textcolor{red}{ \textbf{Cálculo da Energia Consumida por Componente:} utilize a fórmula: $E = V I t$, onde $E$ é a energia em joules, $V$ é a tensão em volts, $I$ a corrente em ampères e $t$ o tempo em segundos. Realize esse cálculo para cada componente do sistema.}
    \item \textcolor{red}{ \textbf{Estimativa do Consumo Total de Energia:} some as energias de todos os componentes para obter o consumo energético total estimado do sistema por lançamento.}
    \item \textcolor{red}{ \textbf{Escolha da Fonte de Alimentação:} com base no consumo total, converta para watt-hora: 1 Wh = 3600 J. Adicione e justifique uma margem de segurança (pesquisar o que é usual para este tipo de sistema) e selecione uma bateria ou fonte que supra a demanda.}
    \item \textcolor{red}{ \textbf{Planejamento do Circuito de Alimentação:} analise a necessidade de inclusão de reguladores de tensão, proteções contra sobrecorrente, e isolamento para atuadores. Como as oscilações de tensão serão controladas?}
    \item \textcolor{red}{ \textbf{Monitoramento via \textit{software}:} implemente, no \textit{software} embarcado, rotinas para medir a tensão da bateria e registrar dados de tempo de operação. Colete os dados para análise.}
    \item \textcolor{red}{ \textbf{Validação com Testes Reais:} realize testes de campo para comparar o consumo estimado com o real. Ajuste a capacidade da fonte se necessário e refine o software de coleta de dados. Justifique a diferença de resultados (real x teórica), de acordo com a literatura.}
\end{itemize}

\section{Descrição de \textit{software}}

\textcolor{red}{Itens fundamentais:}

\begin{itemize}
    \item \textcolor{red}{ \textbf{Diagrama Atividades UML:}  descrever e explicar o fluxo do comportamento funcional do produto proposto, evidenciando, de forma clara e estruturada, como as atividades são executadas, em que ordem e sob quais condições. Esse diagrama permite compreender o funcionamento dinâmico do sistema, destacando:}

    \begin{itemize}
        \item \textcolor{red}{ os principais atores (usuários ou sistemas externos) envolvidos no processo;
        \item as atividades de negócio realizadas por cada ator ou pelo próprio sistema;
        \item os insumos (entradas) necessários para a execução das atividades;
        \item os resultados (saídas) gerados ao longo do fluxo;
        \item os pontos de decisão, paralelismo e sincronização das atividades;
        \item Entre as notações mais relevantes, destacam-se o estado iniciail e final, atividades, nós de decisão e junção, barras de bifurcação, Raias (\textit{swimlanes}) e Fluxos de controle. }
        \end{itemize}

    \item \textcolor{red}{\textbf{\textit{Backlog} do Produto}}
    \begin{itemize}
        \item \textcolor{red}{ Descrever os   requisitos funcionais (RF) com a técnica de documentação e especificação de história de usuário (HU);}
        \item \textcolor{red}{ Descrever os requisitos não-funcionais (RNF) de forma objetiva; Mais facilmente, mais rapidamente, mais responsivo, de fácil uso, são descrições subjetivas não-válidas. Um exemplo de requisito não-funcional de desempenho: ``a página deve carregar em até 5s quando em conexão 4G''. }
        \item \textcolor{red}{ Ao descrever os requisitos (RF e RNF), usar a  classificação MoSCoW (\textit{Must have}, \textit{Should have} e \textit{Could have}),  para auxiliar a priorização dos itens do backlog.}
        \item \textcolor{red}{ Protótipos de interfafe gráfica do \textit{software} em alta fidelidade.}
        \begin{itemize}
            \item \textcolor{red}{\textbf{ Todas HUs devem conter sua descrição (Eu-Como-Para), critérios de aceitação e protótipos de interface, documentadas no github. }}
            \end{itemize} 

        \end{itemize}

    \item \textcolor{red}{ \textbf{Descrição da arquitetura da solução de \textit{software} proposta:}}
    \begin{itemize}
    
        \item \textcolor{red}{ Esta subseção deve contemplar o documento de arquitetura do sistema e deve ser estruturado segundo as visões (4+1) previstas no processo unificado (UP): lógica, de processos;  implementação, implantação e dados (substuirá a visão de casos de uso).}
        
        \item \textcolor{red}{ Propósito do \textit{software} (qual o seu papel no sistema);}
        \item \textcolor{red}{ Padrão adotado: MVC, MVP, Microsserviços, Monolítico, etc (Justificar);}
        \item \textcolor{red}{ Linguagens de programação: Java, Python, C\#, JavaScript, etc.}
        \item \textcolor{red}{ \textit{Frameworks} e bibliotecas: Spring Boot, .NET Core, React, Angular, Django, etc.}
            \item \textcolor{red}{ Banco de dados: Relacional (PostgreSQL, MySQL, etc) X NãoSQL (MongoDB,etc). }
        \item \textcolor{red}{Persistência de dados: Modelo Entidade‑Relacionamento (MER) e seu respectivo Diagrama Entidade‑Relacionamento (DER), aplicáveis quando a solução utiliza banco de dados relacional; alternativamente, diagrama de estrutura de documentos, empregado nos casos em que a arquitetura adota banco de dados não relacional.}
         \end{itemize}     
    
    \item \textcolor{red}{\textbf{ Roteiro de testes funcionais:} }
    \begin{itemize}
        \item \textcolor{red}{Código do caso de teste;}
        \item \textcolor{red}{Nome do caso de teste;}
        \item \textcolor{red}{ Objetivo do caso de teste;}
        \item \textcolor{red}{ Pré-condições do sistema para o teste ser realizado, quando se aplicar;}
        \item \textcolor{red}{ Descrição dos procedimentos a serem executados para o teste;}
        \item \textcolor{red}{ Resultado esperado para o teste ser aprovado (pós-condição após realizado o teste);}
    \end{itemize}
\end{itemize}

